\scanpage{62}
% @ lines 1–5 / text
5) Considérons la forme \deleted{déterminant}\inserted{$\xi\mapsto-\det(\unsure{\vert|\xi})$} sur $M_0=\mathbb M_2(\mathcal O_{S_0})$, et considérons sa restriction au sous-Module $\mathfrak{sl}(2)_{S_0}\subset M_0$ des matrices de trace nulle, \deleted{\unsure{laquelle|le groupe}} i.e. des matrices
% @ line 6 / display
\[
 \xi=\begin{pmatrix}z&x\\y&-z\end{pmatrix}\tag*{(3)}
\]
% @ line 7 / text
donnant
% @ lines 8–9 / display
\[
 \begin{aligned}
 &\deleted{q(\xi)\overset{\mathrm{df}}{=}\det\xi=}\\
 &q_0(\xi)\overset{\mathrm{dfn}}{=}-\det\xi=xy+z^2
 \end{aligned}\tag*{(4)}
\]
% @ lines 10–14 / text
Cette forme \inserted{$-\det u$ sur \deleted{\unsure{$M_0$}} $M_0=\mathbb M_2(\mathcal O_{S_0})$} est évidemment invariante par l'action de $\mathrm{GP}(1)$, déduite par passage au quotient de celle de $\mathrm{Gl}(2)$ -- car
% @ line 15 / display
\[
 \det(g\xi g^{-1})=\det\xi\tag*{(5)}
\]
% @ lines 16–18 / text
a fortiori sa restriction à $\deleted{\unsure{\mathfrak L_0}}\widetilde{\unsure{\mathfrak g|g}}_0=\mathfrak{sl}(2)_{S_0}$ l'est. Identifiant $\widetilde{\mathfrak g}_0$ à $E_0=\mathcal O_{S_0}^3$ par
% @ lines 19–20 / display
\[
 \left\{\begin{aligned}
 (x,y,z)&\longrightarrow\begin{pmatrix}z&x\\y&-z\end{pmatrix}\\
 E_0=\mathcal O_{S_0}^3&\xrightarrow{\sim}\widetilde{\mathfrak g}_0=\mathfrak{sl}(2)_{\unsure{S|S_0}}
 \end{aligned}\right.\tag*{(6)}
\]
% @ lines 21–22 / text
on trouve une action de $G_0=\mathrm{GP}(1)$ sur $(E_0,q_0)$. Il est immédiat
% archive sheet number: 61
